\section{Dataset Preparation}\label{sec:dataset-preparation}

\subsection{ScalarReal Dataset}\label{subsec:scalarreal-dataset}

ScalarReal is this project's local five-video subset of ScalarFlow, introduced by \citet{eckert2019scalarflow} for real-world scalar transport. The public release contains 104 reconstructions, each with 150 frames of 3D density, 3D velocity, rendered images, and input images~\citep{scalarflowData2019}. This subset contains five camera videos and \texttt{info.json}; tables follow JSON field order.

Use labels: \emph{train} is used by Python density training; \emph{test} is held out; \emph{render} builds the synthetic render path; \emph{no} is unused there.

\begin{center}
    \begingroup
    \scriptsize
    \renewcommand{\arraystretch}{1.1}
    \setlength{\tabcolsep}{3pt}
    \begin{tabularx}{\columnwidth}{@{}ll>{\raggedright\arraybackslash}Xccl@{}}
        \toprule
        \textbf{Split} & \textbf{JSON} & \textbf{Files} & \textbf{V} & \textbf{F/view} & \textbf{Use} \\
        \midrule
        train & \texttt{train\_videos} & \texttt{00,01,03,04} & 4 & 120 & train \\
        test & \texttt{test\_videos} & \texttt{02} & 1 & 120 & test \\
        \bottomrule
    \end{tabularx}
    \par\smallskip
    \footnotesize\textbf{Table 1:} ScalarReal split fields.
    \endgroup
\end{center}

\begin{center}
    \begingroup
    \scriptsize
    \renewcommand{\arraystretch}{1.1}
    \setlength{\tabcolsep}{3pt}
    \begin{tabularx}{\columnwidth}{@{}ll>{\raggedright\arraybackslash}Xl@{}}
        \toprule
        \textbf{Field} & \textbf{Value} & \textbf{Meaning} & \textbf{Use} \\
        \midrule
        \texttt{file\_name} & per video & video file path & train \\
        \texttt{frame\_rate} & 30 & video FPS & no \\
        \texttt{frame\_num} & 120 & video length & train \\
        \texttt{camera\_angle\_x} & per video & horizontal FOV & train \\
        \texttt{camera\_hw} & \([1920,1080]\) & image \([H,W]\) & no \\
        \texttt{transform\_matrix} & per video & camera pose & train \\
        \bottomrule
    \end{tabularx}
    \par\smallskip
    \footnotesize\textbf{Table 2:} Per-video fields.
    \endgroup
\end{center}

\begin{center}
    \begingroup
    \scriptsize
    \renewcommand{\arraystretch}{1.1}
    \setlength{\tabcolsep}{3pt}
    \begin{tabularx}{\columnwidth}{@{}ll>{\raggedright\arraybackslash}Xl@{}}
        \toprule
        \textbf{Field} & \textbf{Value} & \textbf{Meaning} & \textbf{Use} \\
        \midrule
        \texttt{frame\_bkg\_color} & \([0,0,0]\) & background color & no \\
        \texttt{voxel\_scale} & \([0.4909,0.73635,0.4909]\) & voxel scale & train \\
        \texttt{voxel\_matrix} & 4x4 & voxel transform & train \\
        \texttt{render\_center} & \([0.338,0.388,-0.261]\) & render-path center & render \\
        \texttt{near} & \(1.1\) & near plane & train \\
        \texttt{far} & \(1.5\) & far plane & train \\
        \texttt{phi} & \(20.0\) & rotation angle & render \\
        \texttt{rot} & \texttt{Y} & rotation axis & render \\
        \bottomrule
    \end{tabularx}
    \par\smallskip
    \footnotesize\textbf{Table 3:} Global fields.
    \endgroup
\end{center}

\subsection{Key Dataset Parameters}\label{subsec:key-dataset-parameters}

This subsection expands only parameters whose role is not already obvious from the tables. The point is the geometry encoded by the dataset, not a particular implementation path.

\subsubsection{Camera Model}\label{subsubsec:dataset-camera-model}

Let \(\theta_x=\texttt{camera\_angle\_x}\), \(h=\texttt{camera\_hw}=[H,W]\), and \(C=\texttt{transform\_matrix}\). The image plane is a pinhole camera grid with horizontal field of view \(\theta_x\). Its focal length is
\[
    f=\frac{0.5W}{\tan(0.5\,\theta_x)}.
\]
The corresponding intrinsic matrix is
\[
    K=
    \begin{bmatrix}
        f & 0 & 0.5W\\
        0 & f & 0.5H\\
        0 & 0 & 1
    \end{bmatrix}.
\]
The pose \(C\) places this camera in the dataset 3D coordinate system:
\[
    C=
    \begin{bmatrix}
        R & o\\
        0 & 1
    \end{bmatrix}.
\]
For image coordinate \((i,j)\), the corresponding camera-space direction is
\[
    d_c=
    \begin{bmatrix}
        (i-0.5W)/f\\
        -(j-0.5H)/f\\
        -1
    \end{bmatrix},
\]
and the dataset-space ray is
\[
    r_o=o,\qquad r_d=R\,d_c.
\]
The negative sign in the second component accounts for image rows increasing downward.

\subsubsection{Volume Transform}\label{subsubsec:dataset-volume-transform}

Let \(M=\texttt{voxel\_matrix}\) and \(s=\texttt{voxel\_scale}\). Split the spatial part of \(M\) into three axis columns and a translation:
\[
    M=
    \begin{bmatrix}
        m_x & m_y & m_z & t\\
        0 & 0 & 0 & 1
    \end{bmatrix}.
\]
For local volume coordinate \(u\in[0,1]^3\), the dataset-space position is
\[
    p(u)=t+u_xs_xm_x+u_ys_ym_y+u_zs_zm_z.
\]
Thus \(M\) supplies orientation and origin, while \(s\) supplies the per-axis metric scale.

\subsubsection{Render-Path Metadata}\label{subsubsec:dataset-render-path}

Let \(c_r=\texttt{render\_center}\), \(z_n=\texttt{near}\), \(z_f=\texttt{far}\), \(\varphi=\texttt{phi}\), and \(a=\texttt{rot}\). The pair \(z_n,z_f\) defines the ray-depth interval attached to the dataset. The same pair also gives the nominal spherical render-path radius
\[
    \rho=\frac{z_n+z_f}{2}.
\]
The center \(c_r\), elevation \(\varphi\), and rotation axis \(a\) describe the synthetic render path. They do not define the volume transform; in particular, \(c_r\) is not the geometric center of the volume.
